\section{Finite factors in their original representations}
\label{sec:finite}

We now apply the boundary construction to close finite factors.  First we
place them on one tracial Hilbert space and identify their two actions of a
hyperfinite subfactor.  This gives the norm control needed for the boundary
argument.  That argument produces a nearby isomorphism; the amplification
estimate then implements it on the original Hilbert space.

\begin{theorem}[Uniform stability of finite factors]
\label{thm:finite}
There are a number $t_{\mathrm{fin}}>0$ and a nondecreasing function
$\eta_{\mathrm{fin}}:(0,t_{\mathrm{fin}})\to(0,\infty)$, with
$\eta_{\mathrm{fin}}(t)\to0$ as $t\downarrow0$, having the following
property.  Let $H$ be a separable nonzero Hilbert space, let
$M\subseteq\B(H)$ be a unital finite factor, and let $N\subseteq\B(H)$
be a von Neumann algebra with the same identity.  If $d(M,N)<t<t_{\mathrm{fin}}$,
then $N$ is a finite factor and there is a unitary $u\in\B(H)$ such that
\[
 uMu^*=N,\qquad \norm{u-I_H}\le\eta_{\mathrm{fin}}(t).
\]
The threshold and function are independent of $M,N$ and their representations.
\end{theorem}

\subsection{An irreducible hyperfinite subfactor}

We use the normalized trace $\tau$ of a finite factor, with
$\norm{x}_2=\tau(x^*x)^{1/2}$ and trace vector
$\Omega=1\in L^2(M,\tau)$.  We use the trace-preserving expectations
onto von Neumann subalgebras of a finite algebra.  An inclusion $R\subseteq M$ is
\emph{irreducible} when $R'\cap M=\C I$.

\begin{lemma}
\label{lem:finite-neighbor}
If $M\subseteq\B(H)$ is a unital finite factor and $N\subseteq\B(H)$
is a von Neumann algebra with $d(M,N)<1/10$, then $N$ is a finite factor.
\end{lemma}
\begin{proof}
Choose $t$ with $d(M,N)<t<1/10$.  If $v\in N$ were a proper isometry,
choose $x\in M_1$ with $\norm{x-v}<t$.  Then
$x^*x\ge(1-t)^2I$, so the polar part $w$ of $x$ is an isometry in $M$.
Finiteness makes $w$ unitary.  Since $1-t\le |x|\le1$,
$\norm{w-v}\le\norm{w-x}+\norm{x-v}<2t<1$.
This is impossible: a proper isometry within distance less than one of a
unitary would be invertible.  Thus $N$ is finite.

Suppose that $q\in Z(N)$ is a projection.  Choose $a\in M_1$ with
$\norm{a-q}<t$.  For $v\in\U(M)$ choose $b\in N_1$ with
$\norm{v-b}<t$.  Since $q$ commutes with $b$,
\[
 \norm{vav^*-q}\le\norm{a-q}+\norm{[v,q]}<3t.
\]
The trace-Hilbert closed convex hull of the conjugates $vav^*$ contains
$\tau(a)I$.  Indeed, its unique vector of least norm is conjugacy
invariant, is bounded by one by weak compactness of the operator ball,
and is therefore scalar; its trace is $\tau(a)$.  Consequently
$\norm{q-\tau(a)I}\le3t<1/2$.  A nontrivial projection has distance
$1/2$ from the scalars, so $q$ is zero or one.  The spectral projections
of every central selfadjoint element are therefore trivial, proving
$Z(N)=\C I$.
\end{proof}

We recall Popa's irreducible hyperfinite subfactor theorem
\cite{Popa1981} and give the finite-factor construction needed here.

\begin{lemma}[Increasing matrix algebras with scalar relative commutant]
\label{lem:irreducible-hyperfinite}
Every separable-predual $\II_1$ factor $M$ contains an irreducible unital
hyperfinite $\II_1$ subfactor $R$.  One can write
$R=(\bigcup_j Q_j)''$, where $Q_j$ are increasing unital full matrix
algebras and their sizes tend to infinity.
\end{lemma}
\begin{proof}
We first make a single selfadjoint element nearly scalar under a matrix
average.  Let $P$ be a $\II_1$ factor, $y=y^*\in P$, and $\eps>0$.
Choose a spectral step approximation $y_0=\sum_{i=1}^l\lambda_i p_i$,
where $\sum_i p_i=I$.  For a sufficiently large integer $b$, split
$\lfloor b\tau(p_i)\rfloor$ orthogonal projections of trace $1/b$ from
each $p_i$.  The remaining projection has trace
\[
 \frac{b-\sum_i\lfloor b\tau(p_i)\rfloor}{b}<\frac lb,
\]
and also splits into projections of trace $1/b$.  The resulting $b$
equivalent projections extend to matrix units for a unital
$F\cong M_b(\C)$ in $P$.  Let $y_1\in F$ be diagonal, equal to
$\lambda_i$ on the pieces taken from $p_i$ and zero on the remainder.
Then
\[
 \norm{y-y_1}_2\le\norm{y-y_0}_2+\norm{y_0}\sqrt{l/b}.
\]
Haar averaging over $\U(F)$ sends $y_1$ to $\tau(y_1)I$.  Contractivity
of this average in $L^2$ gives
\begin{equation}
 \norm{E_{F'\cap P}(y)-\tau(y)I}_2
 \le2\norm{y-y_1}_2<\eps,
 \label{eq:one-element-matrix-average}
\end{equation}
by first choosing $y_0$ and then $b\ge2$.

Choose an $L^2$-dense sequence $(x_j)$ in the selfadjoint unit ball of
$M$, and enumerate the requirements $(j,k)\in\N^2$.  Start with
$Q_0=\C I$.  Given a full matrix algebra $Q_n$, its relative commutant
$P_n=Q_n'\cap M$ is a $\II_1$ factor: matrix units identify $M$ with a
matrix algebra over one of its corners and $P_n$ with that corner.
For the next requirement $(j,k)$ apply
\eqref{eq:one-element-matrix-average} in $P_n$ to $E_{P_n}(x_j)$, with
error $1/k$.  If $F\subseteq P_n$ is the resulting matrix algebra, set
$Q_{n+1}=Q_n\vee F$.  These commuting matrix algebras generate a full
matrix algebra of at least twice the size of $Q_n$, and
\[
 P_{n+1}=F'\cap P_n,\qquad E_{P_{n+1}}=E_{P_{n+1}}E_{P_n}.
\]
Hence $\norm{E_{P_{n+1}}(x_j)-\tau(x_j)I}_2<1/k$, and all earlier
bounds persist under subsequent expectations.  For
$R=(\bigcup_nQ_n)''$ the expectation onto
$R'\cap M=\bigcap_nP_n$ thus sends every $x_j$ to $\tau(x_j)I$.
Density gives $R'\cap M=\C I$.  In particular $R$ is a factor; it is
finite and has unital matrix subalgebras of unbounded size, so it is of
type $\II_1$.  Its increasing matrix presentation makes it hyperfinite.
\end{proof}

We will repeatedly use the following consequence of this construction.
If $E_j$ denotes Haar averaging over $\U(Q_j)$, then $E_j$ is the
orthogonal projection in $L^2(M,\tau)$ onto $L^2(Q_j'\cap M)$.
These decreasing subspaces intersect in $\C I$, so
\begin{equation}
 E_j(x)\longrightarrow\tau(x)I\quad\hbox{in }L^2(M,\tau)
 \qquad(x\in M).
 \label{eq:irreducible-averages}
\end{equation}
Here is a direct justification avoiding any boundedness assumption on a
Hilbert-space vector.  Decreasing orthogonal projections have a strong
Hilbert-space limit.  For bounded $x$, the operators $E_j(x)$ are uniformly
bounded, and every weak-operator cluster belongs to all $Q_j'\cap M$,
hence is scalar; its trace identifies it as $\tau(x)I$.  Applying the
cluster to $\Omega$ identifies the Hilbert-space limit.  Density of the
bounded vectors then gives the assertion on the whole trace Hilbert space.

\subsection{Common tracial position}

A factor $P\subseteq\B(K)$ is in \emph{standard tracial position}
with unit vector $\Omega$ when $x\mapsto x\Omega$ identifies $K$
with $L^2(P,\tau_P)$. Its conjugation $J_P$ is determined by
$J_P(x\Omega)=x^*\Omega$.

The representation theorem we use is the following part of
\citet[Lemma~4.10]{CCSSWW2014}.  We state it separately because it changes
the Hilbert space and does not itself give the unitary in
\Cref{thm:finite}.

\begin{proposition}[Common-position theorem]
\label{prop:finite-common-position}
Let $M,N$ be separable-predual $\II_1$ factors acting nondegenerately on
one Hilbert space, with mutual strict $g$-near inclusions, where
$0<g<1.74\cdot10^{-13}$.  Suppose that an amenable von Neumann subalgebra
$R\subseteq M\cap N$ satisfies $R'\cap M\subseteq R$.  There are
faithful normal representations $\pi:M\to\B(K)$ and
$\rho:N\to\B(K)$, a common unit vector $\Omega$, and a number
$\beta<50948\sqrt g$ such that:
\begin{enumerate}[label=(\roman*),leftmargin=*]
\item both images are in standard tracial position with vector $\Omega$;
\item $\pi|_R=\rho|_R$, and the images have mutual strict
$\beta$-near inclusions;
\item for original contractions $x\in M$, $y\in N$,
\[
 \norm{\pi(x)-\rho(y)}\le\beta+\norm{x-y};
\]
\item if $e_R$ is the projection onto $\overline{\pi(R)\Omega}$,
then $(\pi(M)\cup\{e_R\})''=(\rho(N)\cup\{e_R\})''$.
\end{enumerate}
The space $K$ can be taken separable.
\end{proposition}

Below, after applying the proposition, we suppress $\pi,\rho$
when discussing its images.  Normalizing near-inclusion approximants gives
\begin{equation}
 d(M,N)\le 2\beta.
 \label{eq:standard-gap}
\end{equation}

\begin{lemma}[The common right action]
\label{lem:finite-right-action}
In the setting of \Cref{prop:finite-common-position}, the right actions
of $R$ agree elementwise:
\[
 J_M rJ_M=J_N rJ_N\qquad(r\in R).
\]
In particular $T=J_NJ_M-I$ commutes with the left and right actions of $R$.
\end{lemma}
\begin{proof}
The commutant of $(M\cup\{e_R\})''$ is $J_MRJ_M$.  Indeed, an element
of $M'$ is right multiplication by an element of $M$.  If it commutes
with $e_R$, applying it to $\Omega$ shows that its right multiplier lies
in $L^2(R)\cap M=R$, where the equality follows from the trace
expectation.  Conversely right multiplication by $R$ reduces $L^2(R)$.
The same argument applies to $N$.  Equality of the basic constructions
therefore gives $J_MRJ_M=J_NRJ_N$.  For each $r\in R$ the two displayed
operators take $\Omega$ to $r^*\Omega$.  The vector $\Omega$ is
separating for the common right algebra, since it is cyclic for $M$, so
the operators are equal.  Applying these identities twice proves the last assertion.
\end{proof}

\begin{lemma}[Testing a bimodular operator on bounded algebra elements]
\label{lem:finite-ball-tests}
Let $M$ be a $\II_1$ factor and let $R=(\bigcup_jQ_j)''\subseteq M$
be as in \Cref{lem:irreducible-hyperfinite}.  If a bounded operator $V$
on $L^2(M,\tau)$ commutes with left and right multiplication by $R$, then
\[
 \norm{V}\le384\sup_{a,b\in M_1}
       \abs{\ip{Va\Omega}{b\Omega}}.
\]
\end{lemma}
\begin{proof}
Denote the supremum by $K$.  Start with bounded $a,b\in M$ with
$\norm{a}_2,\norm{b}_2\le1$, without an operator-norm bound.
By \eqref{eq:irreducible-averages}, finite convex averages of conjugations
by unitaries of $R$ can bring both $aa^*$ and $bb^*$ within one of their
scalar traces in $L^2$.  One first takes a sufficiently high Haar average
and then approximates its integral by a finite sum; rational weights and
repetitions give equal weights.  Choose such a family $(d_i)$.
Independently choose a family $(e_l)$ so that the averages of
$e_l^*a^*ae_l$ and $e_l^*b^*be_l$ have the same property.  Enumerate the
Cartesian product of the two families by $k=1,\ldots,m$ and put
$a_k=d_i a e_l$, $b_k=d_i b e_l$.  Their covariance averages satisfy
\[
 L_a=\frac1m\sum_k a_ka_k^*,\quad
 Q_a=\frac1m\sum_k a_k^*a_k,\qquad
 \norm{L_a}_2,\norm{Q_a}_2\le2,
\]
and the same bounds hold for $L_b,Q_b$.

Use the same independent fair signs in the two sums
\[
 A=m^{-1/2}\sum_k\epsilon_k a_k,\qquad
 B=m^{-1/2}\sum_k\epsilon_k b_k.
\]
Bimodularity and sign orthogonality give
\begin{equation}
 \E\ip{VA\Omega}{B\Omega}=\ip{Va\Omega}{b\Omega},
 \qquad \E\norm{A}_2^2,\ \E\norm{B}_2^2\le1.
 \label{eq:finite-random-pairing}
\end{equation}
Put $\norm{x}_4^4=\tau((x^*x)^2)$. The three pairings in the
fourth moment give exactly
\[
 \E\norm{A}_4^4
 =\tau(Q_a^2)+\tau(L_a^2)
  +\frac1{m^2}\sum_{k,l}\tau((a_k^*a_l)^2)
  -\frac2{m^2}\sum_k\tau((a_k^*a_k)^2).
\]
Tracial Cauchy--Schwarz gives
$\abs{\tau((a_k^*a_l)^2)}\le\tau(a_ka_k^*a_la_l^*)$.
Dropping the last nonpositive term therefore yields
\[
 \E\norm{A}_4^4\le\norm{Q_a}_2^2+2\norm{L_a}_2^2\le12,
 \qquad \E\norm{B}_4^4\le12.
\]
If $A=w|A|$, set $A_s=w|A|1_{[0,s]}(|A|)$ and define $B_s$ similarly.
These operators have norm at most $s$, their $L^2$ norms do not increase,
and
\[
 \E\norm{A-A_s}_2^2\le12/s^2,
 \qquad \E\norm{B-B_s}_2^2\le12/s^2.
\]
Substituting them in \eqref{eq:finite-random-pairing} and applying
Cauchy--Schwarz in probability gives
\[
 \abs{\ip{Va\Omega}{b\Omega}}
 \le s^2K+\frac{2\sqrt{12}}s\norm{V}.
\]
Take $s=4\sqrt{12}$, and then take the supremum over the dense bounded
$L^2$ unit vectors $a,b$.  We obtain
$\norm{V}\le192K+\norm{V}/2$, proving the assertion.
\end{proof}

\begin{lemma}[Norm control of tracial conjugations]
\label{lem:tracial-conjugations}
Let $M,N\subseteq\B(K)$ be $\II_1$ factors in common standard tracial
position with vector $\Omega$.  Suppose that $R\subseteq M\cap N$ is
an irreducible hyperfinite subfactor of $M$, presented by increasing
matrix algebras, and that its right actions agree.  If $d(M,N)<b$, then
\[
 \norm{J_M-J_N}\le C_Jb,\qquad C_J=768.
\]
\end{lemma}
\begin{proof}
Set $T=J_NJ_M-I$.  The common right-action identity implies that $T$
commutes with both actions of $R$.  For $x\in M_1$ choose $y\in N_1$
with $\norm{x-y}<b$.  Then
\[
 \norm{Tx\Omega}
 =\norm{J_Nx^*\Omega-x\Omega}
 \le\norm{(x^*-y^*)\Omega}+\norm{(x-y)\Omega}<2b.
\]
Thus the supremum in \Cref{lem:finite-ball-tests} is at most $2b$.
That lemma gives $\norm{T}\le768b$, and
$\norm{T}=\norm{J_M-J_N}$.
\end{proof}

\subsection{From the boundary to a Jordan isomorphism}

We now work entirely in the common tracial representation.  Our immediate
objective, when $d(M,N)<b$ is sufficiently small, is to construct a
Hilbert-space unitary $S_0$ fixing $\Omega$ such that, for every
$u\in\U(M)$,
\[
 S_0u\Omega=b_u\Omega,\qquad b_u\in\U(N),\qquad
 \norm{b_u-u}=O(b),
\]
with the error uniform in $u$.  Unitaries linearly span $M$.  Linearity
of $S_0$ and the separating property of $\Omega$ will make the assignment
$u\mapsto b_u$ extend to a unital linear map from $M$ to $N$.
We will first prove that it preserves the Jordan product
$x\circ y=(xy+yx)/2$; a linear map with this property is called a
\emph{Jordan homomorphism}.  After proving surjectivity, Herstein's
orientation theorem \cite{Herstein1956} will leave only the multiplicative
and antimultiplicative possibilities, and closeness will exclude the
latter.  We record its algebraic proof before constructing $S_0$.
All closeness estimates here are in ordinary operator norm; complete
boundedness enters when we return to the original representation.

\begin{lemma}[The orientation of a Jordan map into a factor]
\label{lem:jordan-orientation}
A surjective linear Jordan homomorphism between unital complex algebras,
whose target is a factor von Neumann algebra, is either a homomorphism
or an antihomomorphism.
\end{lemma}
\begin{proof}
Write the map as $\theta:\mathcal A\to N$, where $N$ is the target factor.
Preservation of the Jordan product gives
preservation of $xyx$, since
$xyx=2x\circ(x\circ y)-x^2\circ y$, where
$x\circ y=(xy+yx)/2$.  Polarizing in $x$ gives preservation of
$xyz+zyx$.  Fix $a,b$, and set
\[
 A=\theta(a),\quad B=\theta(b),\quad D=\theta(ab),\quad
 E=\theta(ba)=AB+BA-D,
 \qquad h=D-AB,\quad k=D-BA.
\]
Compute $\theta(ab\,c\,ba+ba\,c\,ab)$ in two ways.  The polarized
triple identity gives $DCE+ECD$, where $C=\theta(c)$; applying the
unpolarized identity first to the inner and then to the outer products
gives $ABCBA+BACAB$.  Expanding their equality yields
\[
 hCk+kCh=0\qquad(C\in N),
\]
where surjectivity allows every $C$ in the target.  Substitute
$C=C_1kC_2$, and use the same identity with $C_1$ and $C_2$ to move
$h$ past the adjacent factors.  The two resulting terms are equal, so
\[
 hC_1kC_2k=0\qquad(C_1,C_2\in N).
\]
A factor is prime: $xNy=\{0\}$ implies $x=0$ or $y=0$, since the
central supports of the support projections of nonzero $x,y$ are both
one.  If $k\ne0$, primeness first gives $hC_1k=0$ for every $C_1$,
and then $h=0$.  Thus, for each pair $(a,b)$, either
$\theta(ab)=\theta(a)\theta(b)$ or
$\theta(ab)=\theta(b)\theta(a)$.

For fixed $a$, the sets of $b$ satisfying these two respective identities
are additive subgroups whose union is the whole additive group.  A group
cannot be the union of two proper subgroups: choose an element in each
but not the other, and their sum belongs to neither.  One identity
therefore holds for all $b$.  The same two-subgroup argument, now for
those $a$ for which each identity holds for every $b$, makes one
orientation hold globally.
\end{proof}

\begin{proposition}
\label{prop:finite-standard-isomorphism}
There are absolute constants $b_0>0$ and $C_{\mathrm{st}}>0$ such that
the following holds. Let $M,N\subseteq\B(K)$ be type $\II_1$ factors
on a separable Hilbert space in common standard tracial position with
vector $\Omega$. Suppose they contain a common hyperfinite $\II_1$
subfactor $R$, irreducible in $M$ and generated by increasing matrix
algebras, and that $J_MrJ_M=J_NrJ_N$ for every $r\in R$.
If $d(M,N)<b<b_0$, then there is
a normal unital star-isomorphism $\vartheta:M\to N$ satisfying
\[
 \sup_{x\in M_1}\norm{\vartheta(x)-x}\le C_{\mathrm{st}}b.
\]
\end{proposition}
\begin{proof}
Choose a countable strongly dense subgroup $G\subseteq\U(M)$.
For each $g\in G$ choose a unitary $v_g\in N$ with
$\norm{v_g-g}\le3b$: a contraction approximant to $g$ is invertible
for $b<1$, and its polar part has this bound.  Thus
\[
 \sup_{g,h\in G}\norm{v_gv_h-v_{gh}}\le9b.
\]
Apply \Cref{lem:boundary-correction} to these constant fields on the
boundary space $(D,\nu)$ of \Cref{lem:boundary-mean}.  For sufficiently
small absolute $b$ it gives measurable unitary fields $W_g:D\to\U(N)$
satisfying
\begin{equation}
 W_g(x)W_h(g^{-1}x)=W_{gh}(x),\qquad
 \sup_g\norm{W_g-g}_{L^\infty(D)}\le C_1b,
 \label{eq:finite-boundary-cocycle}
\end{equation}
where $C_1$ is absolute.

To obtain this correspondence on unitary vectors, we will compare the
cocycle with the original left and right actions, and then compare its
two boundary values within $N$.  The first comparison uses the separate
action on $D^2$.  Form its product cocycle and the constant representation
\[
 \Sigma_{g,h}(x,y)=W_g(x)J_NW_h(y)J_N,
 \qquad U_{g,h}=gJ_MhJ_M.
\]
The left and right $N$ actions commute, so $\Sigma$ is an exact cocycle
for the separate $G^2$ action.  By
\Cref{lem:tracial-conjugations},
\[
 \sup_{g,h}\norm{\Sigma_{g,h}-U_{g,h}}_\infty
 \le(2C_1+2C_J)b.
\]
The product version of \Cref{lem:boundary-intertwiner} gives a measurable
unitary field $L:D^2\to\U(K)$, with $\norm{L-I}_\infty\le C_2b$, such
that
\begin{equation}
 L(x,y)U_{g,h}=\Sigma_{g,h}(x,y)L(g^{-1}x,h^{-1}y).
 \label{eq:finite-product-intertwiner}
\end{equation}
The constants $C_1,C_2$, and all subsequent $C_i$ in this proof, are
absolute.

The field $L$ takes values in $\U(K)$; its intertwining identity alone
does not supply trace vectors of unitaries in $N$.  For this we use the
diagonal comparison.  Equation~\eqref{eq:finite-boundary-cocycle} gives
\[
 \sup_g\mathop{\rm ess\,sup}_{x,y\in D}
 \norm{W_g(x)W_g(y)^*-I}\le2C_1b.
\]
For sufficiently small $b$, \Cref{lem:boundary-diagonal-unitary}
therefore gives a measurable
$a:D^2\to\U(N)$ with $\norm{a-I}_\infty\le C_3b$ and
\begin{equation}
 a(x,y)=W_g(x)a(g^{-1}x,g^{-1}y)W_g(y)^*.
 \label{eq:finite-diagonal-unitary}
\end{equation}
We can now combine the two comparisons: the vector field
$f(x,y)=L(x,y)^*a(x,y)\Omega$ is equivariant for the diagonal
representation $g\mapsto U_{g,g}$.  Indeed right multiplication
by $W_g(y)^*$ on trace vectors is $J_NW_g(y)J_N$, so
\eqref{eq:finite-diagonal-unitary} and
\eqref{eq:finite-product-intertwiner} give
$f(x,y)=U_{g,g}f(g^{-1}x,g^{-1}y)$.
By \Cref{lem:double-ergodicity}, $f$ is a constant invariant vector.
The invariant vectors for $g\mapsto gJ_MgJ_M$ are exactly $\C\Omega$.
To see this, strong density extends invariance from $G$ to $\U(M)$;
averaging over $\U(Q_j)$ acts on $x\Omega$ as $E_j(x)\Omega$ and
converges to the projection onto $\C\Omega$ by
\eqref{eq:irreducible-averages}.  Hence
\begin{equation}
 L(x,y)^*a(x,y)\Omega=\lambda\Omega,
 \qquad |\lambda|=1,
 \qquad |\lambda-1|\le(C_2+C_3)b.
 \label{eq:finite-boundary-constant}
\end{equation}

Discard one null set and all its countably many $G^2$ translates, so
that the preceding identities hold at a chosen pair $(x,y)$ and at all
its translates.  Define the Hilbert-space unitary
\[
 S_0=\lambda a(x,y)^*L(x,y).
\]
It fixes $\Omega$.  Using \eqref{eq:finite-product-intertwiner} for
$(g,1)$ and \eqref{eq:finite-boundary-constant} at $(g^{-1}x,y)$ gives
\begin{equation}
 S_0g\Omega
   =a(x,y)^*W_g(x)a(g^{-1}x,y)\Omega
   =b_g\Omega,
 \label{eq:finite-unitary-images}
\end{equation}
where $b_g\in\U(N)$ and $\norm{b_g-g}\le C_4b$ uniformly in $g$.

Every $u\in\U(M)$ is the strong limit of a sequence $g_j\in G$.
The vectors $b_{g_j}\Omega=S_0g_j\Omega$ are Cauchy in trace $L^2(N)$.
For unitaries, traciality gives
$\norm{b_{g_j}^*-b_{g_l}^*}_2=\norm{b_{g_j}-b_{g_l}}_2$.
Uniformly bounded $L^2$ convergence in the standard representation
implies strong convergence: test first on the dense right-multiplier
vectors $N'\Omega$ and then use boundedness.  Thus $b_{g_j}$ and their
adjoints converge strongly, and their limits are a unitary $b_u\in N$
and its adjoint.  We obtain
\[
 S_0u\Omega=b_u\Omega,\qquad \norm{b_u-u}\le C_4b,
\]
where the norm inequality follows by weak closedness of the appropriate
operator ball.

Every element of $M$ is a linear combination of unitaries, so the rule
\begin{equation}
 \vartheta(z)\Omega=S_0z\Omega\qquad(z\in M)
 \label{eq:finite-jordan-map}
\end{equation}
defines a linear map into $N$.  It is well-defined because $\Omega$
is separating for $N$.  Writing a contraction as its real and imaginary
selfadjoint parts, each an average of two unitaries, gives
\begin{equation}
 \vartheta(I)=I,\qquad
 \norm{\vartheta(z)-z}\le 2C_4b\norm{z}\quad(z\in M).
 \label{eq:finite-jordan-bound}
\end{equation}
In particular $\vartheta$ is bounded, and it sends every unitary to a
unitary.  For $z=z^*$ expand
$\vartheta(e^{itz})^*\vartheta(e^{itz})=I$ in operator norm.  The
coefficient of $t$ says $\vartheta(z)^*=\vartheta(z)$.  Applying this
also to $z^2$ and comparing the coefficient of $t^2$ gives
$\vartheta(z^2)=\vartheta(z)^2$.  Complex linearity and polarization
show that $\vartheta$ preserves adjoints and the Jordan product.

When $2C_4b<1$, \eqref{eq:finite-jordan-bound} makes its range closed
and makes $\vartheta$ injective.  Every contraction of $N$ has distance
at most $(1+2C_4)b$ from this range, by the original gap and
\eqref{eq:finite-jordan-bound}.  For $(1+2C_4)b<1$ the range must be
all of $N$: otherwise the quotient map by this proper closed linear
subspace has norm one, contrary to that bound.

\Cref{lem:jordan-orientation} shows that this surjective Jordan map is either
multiplicative or antimultiplicative.  The latter possibility is excluded
by \eqref{eq:finite-jordan-bound}.  Indeed, for unitaries $z,w\in M$
antimultiplicativity would give
\[
 \norm{zw-wz}\le\norm{zw-\vartheta(zw)}
    +\norm{\vartheta(w)\vartheta(z)-wz}
 \le6C_4b.
\]
A unital copy of $M_2(\C)$ in $M$ contains anticommuting selfadjoint
unitaries, whose commutator has norm two.  Taking $6C_4b<2$ excludes
the anti case.  Thus $\vartheta$ is a unital star-isomorphism, and it
is normal because a bijective order isomorphism of von Neumann algebras
preserves all bounded increasing suprema.  This proves the proposition,
with $C_{\mathrm{st}}=2C_4$ and one absolute sufficiently small $b_0$.
\end{proof}

\subsection{Return to the original Hilbert space}

\begin{proof}[Proof of \Cref{thm:finite}]
Choose $t> d(M,N)$ below the absolute thresholds specified below.
By \Cref{lem:finite-neighbor}, $N$ is a finite factor.  If either factor
is a matrix algebra, \Cref{thm:amenable} for the mutual near inclusions
gives the required near-identity conjugacy, with an absolute linear
bound.  We may therefore suppose both factors are of type $\II_1$.

Choose $R\subseteq M$ from \Cref{lem:irreducible-hyperfinite}.
By \Cref{thm:amenable}, there is a unitary $w$ with
$\norm{w-I}\le150t$ and $wRw^*\subseteq N$.  Replacing $N$ by
$N_0=w^*Nw$ makes $R$ common, and
\[
 d(M,N_0)<301t.
\]
Take $g=302t$ and apply \Cref{prop:finite-common-position}, provided
$g<1.74\cdot10^{-13}$.  Write $M_s=\pi(M)$, $N_s=\rho(N_0)$.
Their common subfactor remains irreducible in $M_s$ by faithfulness.
By \Cref{lem:finite-right-action} its right actions agree, and
$d(M_s,N_s)\le2\beta$, with $\beta<50948\sqrt g$.
Choose a strict standard-gap bound $b=3\cdot50948\sqrt g$.
If $b<b_0$, \Cref{prop:finite-standard-isomorphism} supplies a normal
star-isomorphism $\vartheta_s:M_s\to N_s$ with
$\norm{\vartheta_s-\id}\le C_{\mathrm{st}}b$ on $M_s$.

Transport it to the original representation:
\[
 \theta=\rho^{-1}\vartheta_s\pi:M\longrightarrow N_0.
\]
For $x\in M_1$ choose $y\in(N_0)_1$ with $\norm{x-y}<g$.
The original-contraction comparison in
\Cref{prop:finite-common-position} and faithfulness of $\rho$ give
\begin{align*}
 \norm{\theta(x)-x}
 &\le\norm{\theta(x)-y}+\norm{y-x}\\
 &\le\norm{\vartheta_s\pi(x)-\pi(x)}
       +\norm{\pi(x)-\rho(y)}+g\\
 &\le C_{\mathrm{st}}b+\beta+2g.
\end{align*}
Let $e(t)=C_{\mathrm{st}}b+50948\sqrt{302t}+604t$.
This is an absolute function tending to zero.

On the original $H\oplus H$, consider the faithful normal paired
representation $x\mapsto x\oplus\theta(x)$ of $M$, and the selfadjoint
switch $s(\xi,\eta)=(\eta,\xi)$.  Its scalar commutator norm on that
factor is at most $e(t)$.  \Cref{thm:amplification} therefore gives
\[
 \norm{\ad_s|_{\{x\oplus\theta(x):x\in M\}}}_{\cb}
 \le C_{\mathrm{amp}}e(t).
\]
At each matrix level the switch commutator has off-diagonal block
$[x_{ij}-\theta(x_{ij})]$, so equivalently
$\norm{\theta-\id}_{\cb}\le C_{\mathrm{amp}}e(t)$.
For sufficiently small absolute $t$, \Cref{lem:paired} gives a unitary
$v\in\B(H)$ with $vMv^*=N_0$ and
$\norm{v-I}\le C_{\mathrm{pair}}C_{\mathrm{amp}}e(t)$, for an absolute
constant $C_{\mathrm{pair}}$.  Then $u=wv$ conjugates $M$ onto $N$, and
\[
 \norm{u-I}\le150t+C_{\mathrm{pair}}C_{\mathrm{amp}}e(t).
\]
All thresholds used above are absolute.  Besides $t<1/10$, they are the
bounds in \Cref{thm:amenable,lem:paired} and the two bounds in
\Cref{prop:finite-common-position,prop:finite-standard-isomorphism}.  Choose $t_{\mathrm{fin}}$ so that they hold for every
$0<t<t_{\mathrm{fin}}$, and take $\eta_{\mathrm{fin}}$ to be the maximum
of the displayed bound and the matrix-factor bound.  These bounds are
nondecreasing functions of $t$ and tend to zero.  This establishes the
claimed uniformity on the original Hilbert space.
\end{proof}
